\documentclass[12pt]{article}
\usepackage[
    ignoreheadfoot,
    top = 1.0 cm,
    right = 1.5 cm,
    left = 1.5 cm,
    bottom = 2.0 cm,
    footskip = 1.0 cm,
]{geometry}
\usepackage{amsmath}
\usepackage{amssymb}
\usepackage{braket}
\usepackage{setspace}
\title{PHYS 285 - HW3}
\author{Kirstin Horvat - 301579693}
\begin{document}
\maketitle
\onehalfspacing


\section*{Question 1a}
1. Find the determinant of A: $det(A - \lambda I)$ \\

\begin{equation}
    \begin{split}
        det(A - \lambda I) &= det\left(\begin{bmatrix} a & b \\ c & d \\ \end{bmatrix} - \lambda \begin{bmatrix} 1 & 0 \\ 0 & 1 \end{bmatrix}\right) \\
        &= det\left(\begin{bmatrix} a - \lambda & b \\ c & d - \lambda \end{bmatrix}\right) \\
        &= (a - \lambda)(d - \lambda) - cb \\
        &= ad - a\lambda - d\lambda + \lambda^2 - cb \\
    \end{split}
\end{equation}
2. Let a = $1$, b = $-a - d$, c = $ad-bc$. Plug into quadratic formula:

\begin{equation}
    \begin{split}
        \lambda_{1,2} &= \frac{-b \pm \sqrt{b^2 - 4ac}}{2a} \\
        &= \frac{-(-a - d) \pm \sqrt{(-a - d)^2 - 4(1)(ad - bc)}}{2} \\
        &= \frac{a + d \pm \sqrt{(a + d)^2 - 4(ad - bc)}}{2} \\\\
        \lambda_1 &= \frac{a + d + \sqrt{(a + d)^2 - 4(ad - bc)}}{2} \\
        \lambda_2 &= \frac{a + d - \sqrt{(a + d)^2 - 4(ad - bc)}}{2} \\
    \end{split}
\end{equation}
\newpage

\text{Eigenvectors of } $A:$
\begin{equation}
\begin{split}
(A - \lambda I)\begin{bmatrix}x \\ y\end{bmatrix} = 0 \\[6pt]
\begin{bmatrix}
a-\lambda & b \\
c & d-\lambda
\end{bmatrix}
\begin{bmatrix}
x \\ y
\end{bmatrix}
=
\begin{bmatrix}
0 \\ 0
\end{bmatrix}
\end{split}
\end{equation}
\begin{equation}
\begin{split}
\lambda_{1,2}
=
\begin{bmatrix}
b \\
\lambda - a
\end{bmatrix}
\quad
\end{split}
\end{equation}


\newpage
\section*{Question 1b}
Eigenvalues for Pauli matrices: \\
\begin{equation}
    \begin{split}
        \sigma_x = \begin{bmatrix} 0 & 1 \\ 1 & 0 \end{bmatrix}, \sigma_y = \begin{bmatrix} 0 & -i \\ i & 0 \end{bmatrix}, \sigma_z = \begin{bmatrix} 1 & 0 \\ 0 & -1 \end{bmatrix} \\\\
    \end{split}
\end{equation}
\begin{equation}
    \begin{split}
        det(\sigma_x - \lambda I) &= det\left(\begin{bmatrix} 0 & 1 \\ 1 & 0 \end{bmatrix} - \lambda \begin{bmatrix} 1 & 0 \\ 0 & 1 \end{bmatrix}\right) \\
        &= det\left(\begin{bmatrix} - \lambda & 1 \\ 1 & - \lambda \end{bmatrix}\right) \\
        &= \lambda^2 - 1\\
        \lambda_1 = 1, \lambda_2 = -1 \\
        \end{split}
\end{equation}

\begin{equation}
    \begin{split}
        det(\sigma_y - \lambda I) &= det\left(\begin{bmatrix} 0 & -i \\ i & 0 \end{bmatrix} - \lambda \begin{bmatrix} 1 & 0 \\ 0 & 1 \end{bmatrix}\right) \\
        &= det\left(\begin{bmatrix} - \lambda & -i \\ i & - \lambda \end{bmatrix}\right) \\
        &= \lambda^2 - 1 \\
        \lambda_1 = 1, \lambda_2 = -1\\
    \end{split}
\end{equation}

\begin{equation}
    \begin{split}
        det(\sigma_z - \lambda I) &= det\left(\begin{bmatrix} 1 & 0 \\ 0 & -1 \end{bmatrix} - \lambda \begin{bmatrix} 1 & 0 \\ 0 & 1 \end{bmatrix}\right) \\
        &= det\left(\begin{bmatrix} 1 - \lambda & 0 \\ 0 & - 1 - \lambda \end{bmatrix}\right) \\
        &= (1 - \lambda)(-1 - \lambda) \\
        \lambda_1 = 1, \lambda_2 = -1
    \end{split}
\end{equation}

\newpage
\textbf{Eigenvectors for Pauli matrices:} \\\\
\text{Eigenvectors of }$\sigma_z$:
\begin{equation}
\begin{split}
\lambda = +1 &: \\
(\sigma_z - I)
\begin{bmatrix}
x \\ y
\end{bmatrix}
=
\begin{bmatrix}
0 & 0 \\
0 & -2
\end{bmatrix}
\begin{bmatrix}
x \\ y
\end{bmatrix}
= 0 \\
 y = 0 \\
\ket{\uparrow}
=
\begin{bmatrix}
1 \\ 0
\end{bmatrix}
\\[6pt]
\lambda = -1 &: \\
(\sigma_z + I)
\begin{bmatrix}
x \\ y
\end{bmatrix}
=
\begin{bmatrix}
2 & 0 \\
0 & 0
\end{bmatrix}
\begin{bmatrix}
x \\ y
\end{bmatrix}
= 0 \\
x = 0 \\
\ket{\downarrow}
=
\begin{bmatrix}
0 \\ 1
\end{bmatrix}
\end{split}
\end{equation}

\text{Eigenvectors of } $\sigma_x$:
\begin{equation}
\begin{split}
\lambda = +1 &: \\
(\sigma_x - I)
\begin{bmatrix}
x \\ y
\end{bmatrix}
=
\begin{bmatrix}
-1 & 1 \\
1 & -1
\end{bmatrix}
\begin{bmatrix}
x \\ y
\end{bmatrix}
= 0 \\
x = y \\
\ket{+}
=
\frac{1}{\sqrt{2}}
\begin{bmatrix}
1 \\ 1
\end{bmatrix}
\\[6pt]
\lambda = -1 &: \\
(\sigma_x + I)
\begin{bmatrix}
x \\ y
\end{bmatrix}
=
\begin{bmatrix}
1 & 1 \\
1 & 1
\end{bmatrix}
\begin{bmatrix}
x \\ y
\end{bmatrix}
= 0 \\
x = -y \\
\ket{-}
=
\frac{1}{\sqrt{2}}
\begin{bmatrix}
1 \\ -1
\end{bmatrix}
\end{split}
\end{equation}


\text{Eigenvectors of } $\sigma_y$:
\begin{equation}
\begin{split}
\lambda = +1 &: \\
(\sigma_y - I)
\begin{bmatrix}
x \\ y
\end{bmatrix}
=
\begin{bmatrix}
-1 & -i \\
i & -1
\end{bmatrix}
\begin{bmatrix}
x \\ y
\end{bmatrix}
= 0 \\
y = i x \\
= \frac{1}{\sqrt{2}}
\begin{bmatrix}
1 \\ i
\end{bmatrix}
\\[6pt]
\lambda = -1 &: \\
(\sigma_y + I)
\begin{bmatrix}
x \\ y
\end{bmatrix}
=
\begin{bmatrix}
1 & -i \\
i & 1
\end{bmatrix}
\begin{bmatrix}
x \\ y
\end{bmatrix}
= 0 \\y = -i x \\
= \frac{1}{\sqrt{2}}
\begin{bmatrix}
1 \\ -i
\end{bmatrix}
\end{split}
\end{equation}

\newpage
\section*{Question 2}

Uncertainty: $\Delta\sigma^2 = \braket{\sigma^2} - \braket{\sigma}^2 $ \\
We also know: $S_x = S_y = S_z = \frac{\hbar}{2}$ \\

\begin{equation}
    \begin{split}
        \ket{\Psi} = \frac{3}{5}\ket{\uparrow} + \frac{4}{5}\ket{\downarrow} = \frac{3}{5}\begin{bmatrix} 1 \\ 0 \end{bmatrix} + \frac{4}{5}\begin{bmatrix} 0 \\ 1 \end{bmatrix} = \begin{bmatrix} \frac{3}{5} \\ \frac{4}{5} \end{bmatrix} \\
            \end{split}
\end{equation}

\begin{equation}
    \begin{split}
            \braket{S_x} &= \braket{\Psi|S_x|\Psi} \\
            &= \frac{\hbar}{2}\begin{bmatrix} \frac{3}{5} & \frac{4}{5} \\ \end{bmatrix} \begin{bmatrix} 0 & 1 \\ 1 & 0 \end{bmatrix} \begin{bmatrix} \frac{3}{5} \\ \frac{4}{5} \end{bmatrix} \\
            &= \frac{\hbar}{2}\left[\frac{3}{5}\cdot\frac{4}{5} + \frac{4}{5}\cdot\frac{3}{5}\right] = \frac{\hbar}{2}\cdot\frac{24}{25} = \frac{12\hbar}{50}
    \end{split}
\end{equation}
\begin{equation}
    \begin{split}
        \Delta S_x^2 &= \braket{S_x^2} - \braket{S_x}^2 \\
        &= \sqrt{\frac{\hbar^2}{4} - \frac{144\hbar^2}{625}} \\
        &= \frac{7}{50}\hbar \\
    \end{split}
\end{equation}

\begin{equation}
    \begin{split}
            \braket{S_y} &= \braket{\Psi|S_y|\Psi} \\
            &= \frac{\hbar}{2}\begin{bmatrix} \frac{3}{5} & \frac{4}{5} \\ \end{bmatrix} \begin{bmatrix} 0 & i \\ -i & 0 \end{bmatrix} \begin{bmatrix} \frac{3}{5} \\ \frac{4}{5} \end{bmatrix} \\
            &= \frac{\hbar}{2}\left[\frac{3}{5}\cdot\frac{4i}{5} - \frac{4}{5}\cdot\frac{3i}{5}\right] = 0 \\
    \end{split}
\end{equation}
\begin{equation}
    \begin{split}
        \Delta S_y^2 &= \braket{S_y^2} - \braket{S_y}^2 \\
        &= \sqrt{\frac{\hbar^2}{4} - 0} \\
        &= \frac{\hbar}{2} \\
    \end{split}
\end{equation}

\begin{equation}
    \begin{split}
            \braket{S_z} &= \braket{\Psi|S_z|\Psi} \\
            &= \frac{\hbar}{2}\begin{bmatrix} \frac{3}{5} & \frac{4}{5} \\ \end{bmatrix} \begin{bmatrix} 1 & 0 \\ 0 & -1 \end{bmatrix} \begin{bmatrix} \frac{3}{5} \\ \frac{4}{5} \end{bmatrix} \\
            &= \frac{\hbar}{2}\left[\frac{3}{5}\cdot\frac{3}{5} + \frac{4}{5}\cdot\frac{4}{5}\right] = \frac{\hbar}{2}\frac{-7}{25} = \frac{-7\hbar}{50}
    \end{split}
\end{equation}
\begin{equation}
    \begin{split}
        \Delta S_z^2 &= \braket{S_z^2} - \braket{S_z}^2 \\
        &= \sqrt{\frac{\hbar^2}{4} - \frac{49\hbar^2}{2500}} \\
        &= \frac{24}{50}\hbar \\
        &= \frac{12}{25}\hbar \\
    \end{split}
\end{equation}

\newpage
\section*{Question 3a}

\begin{equation}
    \begin{split}
        \ket{\Psi(t)} &= e^{-iHt/\hbar}\ket{\Psi(0)} = e^{-iHt\hbar}\ket{\uparrow} \\
        &= e^{-iE_0t/\hbar}\ket{\uparrow} \\
    \end{split}
\end{equation}
\begin{equation}
    \begin{split}
        P_\uparrow &= \left|\braket{\uparrow|\Psi(t)}\right|^2 \\
        &= \left| \bra{\uparrow} \left(e^{-iE_0t/\hbar}\ket{\uparrow}\right) \right|^2 \\
        &= \left|e^{-iE_0t/\hbar}\right|^2 \\
        &= 1 \\
\end{split}
\end{equation}
\text{Global phase factor does not effect probability, so probability is 1} 

\section*{Question 3b}

\begin{equation}
    \begin{split}
        \ket{\Psi(t)} &= e^{-iHt/\hbar}\ket{\Psi(0)} \\
        &= e^{-iHt\hbar}\left[ cos(\frac{\theta}{2})\ket{\uparrow} + e^{i\phi}sin(\frac{\theta}{2})\ket{\downarrow}\right] \\
        &= e^{-iE_0t/\hbar}cos(\frac{\theta}{2})\ket{\uparrow} + e^{-iE_1t/\hbar} e^{i\phi}sin(\frac{\theta}{2})\ket{\downarrow} \\
        &=  e^{-iE_0t/\hbar}\left[cos(\frac{\theta}{2})\ket{\uparrow}+ e^{i(\phi - \frac{(E_1 - E_0)t}{\hbar})}sin(\frac{\theta}{2})\ket{\downarrow}\right]
    \end{split}
\end{equation}

\begin{equation}
    \begin{split}
        P_\uparrow &= \left|\braket{\uparrow|\Psi(t)}\right|^2 \\
        &= \left| \bra{\uparrow}\left[ e^{-iE_0t/\hbar}cos(\frac{\theta}{2})\ket{\uparrow} + e^{-iE_1t/\hbar} e^{i\phi}sin(\frac{\theta}{2})\ket{\downarrow}\right]\right|^2 \\
        &= \left|cos(\frac{\theta}{2})\right|^2 \\
        &= cos^2(\frac{\theta}{2}) \\
\end{split}
\end{equation}

\begin{equation}
    \begin{split}
        P_+ &= \left|\braket{+|\Psi(t)}\right|^2 \\
        &= \left| \left[ \frac{1}{\sqrt{2}}\bra{\uparrow} + \frac{1}{\sqrt{2}}\bra{\downarrow}\right] \left[ e^{-iE_0t/\hbar}cos(\frac{\theta}{2})\ket{\uparrow} + e^{-iE_1t/\hbar} e^{i\phi}sin(\frac{\theta}{2})\ket{\downarrow}\right]\right|^2 \\
        &= \left|e^{-iE_0t/\hbar}\left[\frac{1}{\sqrt{2}}cos(\frac{\theta}{2}) + \frac{1}{\sqrt{2}}e^{i(\phi - \frac{(E_1 - E_0)t}{\hbar})}sin(\frac{\theta}{2})\right]\right|^2 \\
        &= \frac{1}{2}\left| \cos\left(\frac{\theta}{2}\right) + e^{i(\phi - \frac{(E_1 - E_0)t}{\hbar})} \sin\left(\frac{\theta}{2}\right) \right|^2 \\
        &= \frac{1}{2} \left[ \cos^2\left(\frac{\theta}{2}\right)  + \sin^2\left(\frac{\theta}{2}\right) + 2\cos\left(\frac{\theta}{2}\right) \sin\left(\frac{\theta}{2}\right) \cos\left(\phi - \frac{(E_1 - E_0)t}{\hbar}\right) \right] \\\\
        &= \frac{1}{2} \left[ 1 + \sin\theta \cos\left(\phi - \frac{(E_1 - E_0)t}{\hbar}\right) \right] \\
    \end{split}
\end{equation}

\end{document}
